\chapter{Environment}
\section{Simulation}
This project is developed as part of an academic simulation within the Internet of Things course, and is therefore designed to operate within the controlled environment and assumptions defined by the professors, which are outlined briefly below.

\section{Contiki-NG}
\textbf{Contiki-NG} \cite{contiki_ng} represents a sophisticated, open-source operating system explicitly architected for resource-constrained embedded devices within the Internet of Things ecosystem. 

At the foundation of its architecture lies a highly modular, event-driven kernel that operates utilizing Protothreads. This specialized programming model facilitates memory-efficient cooperative multithreading without incurring the heavy overhead typically associated with per-thread stacks, thereby allowing the entire operating system to function optimally within remarkably tight hardware constraints, often requiring as little as 10 kilobytes of random-access memory and a total code footprint of approximately 100 kilobytes.

\section{Cooja}
\textbf{Cooja} is an advanced, Java-based open-source network simulator inherently integrated within the Contiki and Contiki-NG ecosystems, specifically engineered to emulate the intricate dynamics of IoT and wireless sensor networks. It supports cross-level simulation, allowing analysis from network topology down to operating system behavior and machine code execution. It can emulate hardware using tools like MSPSim, enabling real embedded code to run as it would on devices such as MSP430-based nodes.

Cooja offers different mote types for testing: fast "Cooja motes" for high-level simulation and emulated motes for detailed hardware-level validation, including drivers and power management. With a graphical interface for visualizing networks, timelines, and serial output, it helps developers analyze radio interactions, debug protocols, and optimize systems before deploying on real hardware.

\section{Tunslip}
\textbf{Tunslip} is a utility included in Contiki-NG that establishes a virtual SLIP (Serial Line Internet Protocol) link between a border router and a host computer. It is primarily used to bridge an IoT network with the host's networking stack by creating a tunnel interface.

Through this virtual connection, IPv6 packets generated within a simulated or physical sensor network can be forwarded to and from external networks, including the internet. This makes it possible for constrained IoT nodes to communicate with standard IP-based systems as if they were part of the same routed network, facilitating testing and integration of networked applications.

\section{Nordic Dongle NRF52840}
The \textbf{nRF52840 Dongle} is a compact, low-cost USB development platform from \textit{Nordic Semiconductor} built around the nRF52840 SoC. It uses a 32-bit ARM Cortex-M4 processor with a hardware FPU running at 64 MHz and includes 1 MB of Flash and 256 KB of RAM for embedded applications.

It supports multiple 2.4 GHz wireless protocols, including Bluetooth 5 (BLE, long range, 2 Mbps, and mesh), IEEE 802.15.4 for Thread and Zigbee, as well as ANT and proprietary protocols. Security is handled by the integrated CryptoCell-310 hardware accelerator for cryptographic operations.

The board is USB-powered, operates from 1.7 V to 5.5 V, and includes a user button, LEDs, and 15 GPIOs accessible via edge pads. It can be programmed and updated via USB and integrates with the nRF Connect for Desktop ecosystem without requiring an external debugger.
